\section{Related Work}
\label{appx:related_work}

\subsection{NNGP}
\label{appx:nngp}

An hierarchical Bayesian approach to nearest neighbour GPs is derived in \cite{Datta2016}, who construct a full stochastic process allowing an end-to-end probabilistic approach they term `NNGP' derived from a `parent' GP using collections of nearest neighbour sets forming a `reference' set. This work is modified in various ways in \cite{Finley2019} including by adaption to a hybrid empirical-Bayes/fully-Bayesian approach to improve scalability. This is described in Algorithm 5 which presents an MCMC-free approach (`conjugate NNGP') where some parameters are estimated via K-fold cross-validation and the remainder are given conjugate priors to allow exact posterior inference. Under this model, the marginal predictive distribution is Student-\(t\) with mean and variance expressions that match that given by GPnn up to a scaling factor on the variance. Here we will go into more detail to make this connection explicit. First of all, note that they use an alternative parameterisation to us, specifying an inverse lengthscale (\(\phi\)), a ratio (\(\alpha\)) of noise variance (\(\tau^2\)) to kernelscale (\(\sigma^2\)) and a vector of coefficients for their linear mean function (\(\bm{\beta}\), which we can neglect since we focus exclusively on mean-zero GPs). 

In Algorithm 5 they go on to give the following expressions for the predictive mean and variance:
\begin{align*}
    z &= M(s,N(s,k)) \\
    w &= \solve(M[N(s,k),N(s,k)],z) \\
    m_0 &= \widehat{y(s)} = \textrm{dot}(x(s), g(k)) + \textrm{dot}(w,(y[N(s,k)] - \textrm{dot}(X[N(s,k),],g(k)))) \\
    u &= x(s) - \textrm{dot}(X[N(s,k),],w) \\
    v_0 &= \textrm{dot}(u, \textrm{gemv}(V(k),u)) + 1 + \alpha - \textrm{dot}(w,z)  \\
    \widehat{\Var(y(s))} &= b^*_\sigma(k)v_0/(a^*_\sigma(k) - 1) 
\end{align*}

We can see that the predictive mean \(m_0\) matches our own in the mean-zero setting. In particular, in this case the only non-zero term is \(\textrm{dot}(w, y[N(s,k)])\) which can be re-written in our notation as \(\kstar_N^TK_N^{-1}\bm{y}_N\). Note that this expression is, for fixed \(\alpha\), independent of choice of kernelscale. This can be seen using the notation given in \autoref{intro}: \(\kstar_N^TK_N^{-1}\bm{y}_N = \sigma_f^2\cstar_N^T(\sigma_f^2 C_N + \sigma_\xi^2 I)^{-1}\bm{y}_N = \cstar_N^T(C_N+\alpha I)^{-1}\bm{y}_N\).

In addition to this, the predictive variance is equivalent to ours up to multiplicative scaling \( S = b^*_\sigma(k)/[\sigma_f^2 (a^*_\sigma(k) - 1)] \); i.e. $\widehat{\Var(y(s))} = S \predvar_N$, which be seen as follows: In the mean-zero case we have \(v_0 = 1 + \alpha - \textrm{dot}(w,z)\) which can again be re-written in our notation as \(1 + \sigma_\xi^2/\sigma_f^2 - \cstar_N^TC_{\alpha,N}^{-1}\cstar_N\) (where \(C_{\alpha,N}=C_N+\alpha I\)). Hence, if we rescale this by a factor of \(\sigma_f^2\) we exactly obtain the GPnn predictive variance, so that  $\widehat{\Var(y(s))} = b^*_\sigma(k)v_0/(a^*_\sigma(k) - 1) = b^*_\sigma(k)\predvar_N /[\sigma_f^2 (a^*_\sigma(k) - 1)]$ as claimed. Since GPnn uses an additional recalibration step to rescale all predictive variances by a single shared multiplicative factor, the factor $S$ between the two methods can be made effectively redundant: i.e. if our recalibration step (\autoref{alg:gpnn_calibration}) were applied to their method, the two methods would become equivalent {\it with respect to RMSE and weak-calibration} (setting aside differences in parameter estimation and their associated costs).

We reemphasise that whilst GPnn and Conjugate NNGP pointwise predictions can be related as above, the latter is given in the form of a \(t\)-distribution which, even with matched first two moments, will have a different shape to that of a Normal distribution. As such, the (Gaussian) NLL performance measure is no longer a valid choice to assess this model. RMSE and weak-calibration remain distribution-agnostic however.